\chapter{The H$_2$ Covalent Bond and the Helium Limit}\label{ch:h2}
\index{H2@H$_2$ molecule}\index{helium atom}\index{covalent bond}\index{Kolos--Wolniewicz reference}\index{R\"udenberg--Kutzelnigg}\index{bond, helium limit}

% Material to be added later from prel book (realquantum2.tex):
%   - chapters/H2molecule.tex (H2: bonded and non-bonded — 2017 treatment)

\begin{quote}
\itshape It is nonsense to talk about the trajectory of an electron inside an atom.\\
\normalfont\hfill --- E.~Schr\"odinger to M.~Born, 5th Solvay Conference, Brussels, 1927.
\end{quote}
\bigskip

The hydrogen molecule H$_2$ is the simplest covalent bond: two protons sharing two electrons. It is the natural first test of how covalent bonding emerges from the multiphase formulation, for three reasons. First, it can be done parameter-free at Level 1 --- there is no kernel softening to calibrate, no architecture choice beyond the symmetric two-electron partition, no fitted constant. Second, the answer is known to high precision from the Kolos--Wolniewicz reference calculation~\cite{Kolos1968}, against which any new method is measured. Third, the H$_2$ bond connects continuously, through a one-parameter family in the proton separation $R$, to the Helium atom at $R = 0$: \emph{the covalent bond and the atomic shell are the same minimum-energy electron-packing principle at different nuclear separations}.

This chapter works through the H$_2$ result, identifies the mechanism by which the Bernoulli condition makes covalent bonding possible, and follows the $R \to 0$ continuation into the Helium limit.

\section{Parameter-free H$_2$ at Level 1}\label{sec:h2-level1}

H$_2$ at Level 1 is two unit electron densities that do not overlap (Bernoulli partition), distributed around two H$^+$ point kernels at separation $R$. The two electrons occupy opposite halves of the inter-nuclear region --- the natural minimum-energy partition for two electrons in a symmetric two-kernel field. Each electron is attracted to both nuclei and repelled by the other electron through Coulomb interactions only. There is no exchange term, no antisymmetric wave function on $\mathbb{R}^6$, no orbital basis. The setup is the H$_2$ analogue of Chapter~\ref{ch:atoms}'s atomic model: nuclei, electrons, Coulomb energy, non-overlap constraint, variational minimum.

At $R = 1.40$~a.u.\ the converged parameter-free Level-1 calculation gives a total energy
\[
E_{\rm RealQM}(1.40) \;=\; -1.2750 \text{ Ha},
\]
to be compared with the Kolos--Wolniewicz exact value
\[
E_{\rm KW}(1.40) \;=\; -1.174476 \text{ Ha},
\]
an over-binding of $8.6\%$ --- an ordinary total energy for this model, in the same range as
helium at $4.2\%$ and the atomic totals of Chapter~\ref{ch:atoms} at $1.6$--$4.5\%$. The
corresponding binding energy is
\[
\Delta E_{\rm bind} \;=\; 2 E_{\rm H} - E_{\rm H_2}(1.40) \;=\; 0.2747 \text{ Ha}
\;=\; 172 \text{ kcal/mol},
\]
against the measured dissociation energy of $109.5$~kcal/mol: \textbf{$57\%$ too high}.

\paragraph{Why these two figures differ so much, and which one to read.}
Nothing extra has gone wrong at the bond. The binding energy is $E_{\rm bind}/E_{\rm total} =
0.1745/1.1745 = 14.9\%$ of the total, so a per-cent error on the total is necessarily amplified by
its reciprocal when the difference is formed: $8.6\%/0.149 = 57\%$, which is the whole of the
discrepancy. The total is the figure that characterises the model; the binding energy is a
difference one seventh the size of the quantity whose error it inherits. The amplification is the
same structural fact that governs the ionization energies of Chapter~\ref{ch:atoms}, and it is why
``$+57\%$'' must never be quoted without it --- read alone it invites the conclusion that the model
cannot bond, which is not what the number says.

The error also fails to cancel in this particular difference, and for a specific reason: two
separated H atoms carry no shared electron interface, while H$_2$ carries one. The interface error
therefore lands on one side of the subtraction only. Differences taken within a \emph{fixed}
composition --- the same atoms rearranged --- do cancel it, which is why geometries and forces come
out so much better than energies.

\paragraph{An earlier figure withdrawn.} Previous drafts of this chapter reported
$E \approx -1.18$~Ha, $\sim\!0.5\%$ against Kolos--Wolniewicz, and a binding energy of
$-112$~kcal/mol within $3\%$ of measurement. Those values came from runs stopped before
convergence. The relaxation passes through the exact value on its way to the construction's own
minimum, so an under-converged run can land near $-1.17$ by accident, and three of them did. The
figures above are converged on a five-rung mesh ladder with the kernel positions held commensurate
with the grid, and they are what this construction actually gives.

\paragraph{The bond length, and the artefact that hid it.}
The minimum is genuinely shallow --- the energy varies by about $1$~kcal/mol over
$R \in [1.4, 1.6]$~a.u. --- but shallowness was not what obstructed the determination of $R_{\rm
eq}$. The obstruction was grid commensuration. A point kernel sits at $\pm R/2$ on a fixed mesh, so
as $R$ varies the kernels move \emph{relative to the cells} and the discretisation error
oscillates. The amplitude of that oscillation is about $0.07$~Ha/bohr, the same size as the force
itself near equilibrium, while the energy differences between neighbouring $R$ are smaller still.
Every bond length read off such a scan lies inside the noise, and the apparent non-monotonicity of
the force --- which no real diatomic curve shows --- was its signature.

The fix is specific and cheap: scan only separations commensurate with the mesh, so the kernels land
exactly on cell centres and the error is constant instead of oscillating. The force is then monotone
in $R$ at every mesh, and the two independent determinations of the minimum agree ---
$R_{\rm min} = 1.3682$ from the energy and $1.3702$ from the vanishing of the electrostatic force,
$0.15\%$ apart. Against the measured $1.401$~a.u.:
\[
R_{\rm eq}^{\rm RealQM} \;=\; 1.368 \text{ a.u.}, \qquad \textbf{2.3\% short},
\]
stable across meshes. The force constant comes out $0.37$--$0.44$~Ha/bohr$^2$ against the exact
$0.370$: it brackets the right value but moved $19\%$ between meshes, so it is not converged and
the agreement should not be quoted as a figure. \emph{The general lesson applies to every
derivative with respect to a nuclear position:} a point kernel on a fixed mesh contaminates all of
them, and the scan must be commensurate before any geometry is read from it.

So the geometry is right to a few per cent and the energy is over-bound by a few per cent, the two
errors being unrelated in size because the second is amplified by the differencing.

\paragraph{A note on the variational comparison.}
RealQM and the Kolos--Wolniewicz treatment solve different mathematical problems: Kolos--Wolniewicz minimises an antisymmetric two-electron wave function on $\mathbb{R}^6$ over the standard Coulomb Hamiltonian, while RealQM minimises a sum of two unit-density wave functions on non-overlapping spatial supports in $\mathbb{R}^3$ over the multiphase variational principle (Chapter~\ref{ch:equation}). The two are different mathematical objects and the standard ``variational lower bound'' relation between them does not strictly apply --- one is not a restricted variational space of the other. The numerical comparison above is therefore between two independently-computed predictions for the same physical system, not a statement that one bounds the other. This also disposes of an objection that would otherwise be fatal: RealQM's split construction returns $-1.2750$, \emph{below} the exact $-1.174476$, and a variational approximation to the Schr\"odinger Hamiltonian cannot go below the true ground state. RealQM is not one --- different functional, no antisymmetry, per-domain kinetic energy --- so going below is what a different functional does, not a numerical failure.

The bond emerges from the same variational principle that organises the atomic shells of Chapter~\ref{ch:atoms}: the electrons reorganise their territories to minimise the total Coulomb energy, and at the equilibrium $R$ the two-kernel two-electron configuration is energetically preferred over two isolated H atoms. There is no separate ``bond'' machinery; the bond is what the variational principle does when more than one nucleus is present.

\section{Two constructions at $N=2$, and a bracket on the exact energy}\label{sec:h2-bracket}

At two electrons --- and only there --- the multiphase construction can be closed in two different
ways. Both are worth computing, because they fail in opposite directions and together they say
something the single number above cannot.

\paragraph{Split: the non-overlapping form.} The construction used throughout this book. The two
electrons occupy the half-spaces either side of the plane bisecting the internuclear axis, meeting
at a free boundary with the Bernoulli condition of Chapter~\ref{ch:equation}. Each electron feels
the other's potential and not its own, so the pair is counted once and self-repulsion is excluded
by construction.

\paragraph{Hom: the overlapping form.} Both electrons carried by the \emph{same} density over the
whole domain --- sitting on top of each other. Non-overlap is then abandoned, and something must
replace it, because a single blob of charge $2$ and density $2\rho$ has Coulomb self-energy
$\tfrac12\!\int\!\!\int(2\rho)(2\rho)/|x-y| = 2J$ where only one pair is present. Reducing the
repulsion energy by a factor two restores the correct count, giving
\begin{equation}
E \;=\; 2\,(T + V_K) \;+\; J, \qquad
\bigl(-\tfrac12\nabla^2 - K + P[\rho]\bigr)u \;=\; \varepsilon u ,
\label{eq:hom2}
\end{equation}
each electron seeing the other's unit charge through the potential $P[\rho]$.

Both were computed in one code on one five-rung mesh ladder, by axisymmetric reduction --- which is
exact here, since the domains and the kernel potential are all invariant about the internuclear
axis, so nothing is discarded in passing to two dimensions. Two controls fix what the
discretisation contributes. One electron on the same meshes gives $-0.500179$ for H and
$-2.001615$ for He$^+$, errors of $0.04\%$ and $0.08\%$. And \eqref{eq:hom2} has a known answer to
be checked against: it is the Hartree functional, which for a closed two-electron shell coincides
with Hartree--Fock. The Hartree--Fock limits must therefore come out, and they do --- $-2.86111$
against $-2.861680$ for He, and $-1.13363$ against $-1.133630$ for H$_2$, to $0.02\%$ and to all
five figures. The errors below are the models' and not the mesh's.

The exact energies lie between the two constructions:
\begin{align*}
\text{He:}\quad & \underbrace{-2.8611}_{\text{hom},\;+1.5\%} \;<\;
                  \underbrace{-2.903724}_{\text{exact}} \;<\;
                  \underbrace{-3.0265}_{\text{split},\;-4.2\%} \\[2pt]
\text{H}_2\ (R=1.4):\quad & \underbrace{-1.1336}_{\text{hom},\;+3.5\%} \;<\;
                  \underbrace{-1.174476}_{\text{exact}} \;<\;
                  \underbrace{-1.2750}_{\text{split},\;-8.6\%}
\end{align*}
and the same bracket holds on the binding energy, more widely because the binding is a seventh of
the total: $0.1333$~Ha for hom and $0.2747$ for split against the exact $0.1745$, which is $24\%$
low and $57\%$ high.

Overlapping the densities completely and halving the repulsion \emph{under}-binds; separating them
completely \emph{over}-binds. The truth sits between, which is what one would expect of two
idealisations of the same interface: a shared density gives no mutual avoidance beyond the mean
field, disjoint domains give complete and permanent avoidance, and the real state avoids partially.
Non-overlap is in effect a dial between no avoidance and total avoidance, and the exact energies of
both systems lie inside the range the model's own two treatments span. That is a more informative
account of the discrepancy than ``the model over-binds''.

\paragraph{Why split is the model.}
Only one of the two generalises, and this decides which is the model and which is a diagnostic.
Beyond $N=2$ the hom construction has no definite content: putting all $N$ electrons into one
density requires a rule for what replaces non-overlap, and no choice is canonical. A charge-$N$
blob with the repulsion scaled by $(N-1)/N$ is one guess, and it has no shell structure whatever
--- every electron sits in the same lowest state, so there is no atom past helium and no periodic
table. Restoring distinct states requires them to be mutually orthogonal, which is antisymmetry,
which is the $3N$-dimensional formulation again. The split construction needs no such decision at
any $N$: $N$ densities on $N$ non-overlapping domains, every pair counted once by the $i \neq j$
rule, exactly as at $N=2$. \emph{Hom is a bracketing device for two electrons; split is the model
for all $N$.}

\paragraph{Why two electrons are representative.}
The few per cent found here is not an artefact of a small system; it is the accuracy the model has
everywhere, and that is the point on which Chapter~\ref{ch:precision}'s comparison rests. Across
the first three periods the mean error on atomic totals is $1$--$2\%$ with a worst case of $4.5\%$;
He and H$_2$ sit at $1.5$--$8.6\%$ depending on the construction; alkali ionization energies come
within $1.4\%$; hydrogen-bond geometries and the secondary structure of a $153$-residue polypeptide
are right to a few per cent. One electron, two, eighteen, ten thousand --- the figure does not move
much, because its source does not change with $N$. The error lives in the treatment of the
interface between two charge densities, and a system of ten thousand electrons has the same kind of
interface as a system of two, merely more of them. Adding electrons multiplies the boundaries
without altering what happens at any one of them.

\section{Bonding mechanism: density accumulation at the free boundary}\label{sec:bonding-mechanism}

The bond is established by \textbf{accumulation of electron density between the two kernels}, which lowers the total kernel-attraction potential energy by bringing density closer to both nuclei. In the unit-density formulation, the two non-overlapping electron territories meet at the free boundary between the kernels, where \emph{both densities are non-zero}: neither is required to vanish at the interface.

This non-vanishing meeting is precisely what makes the inter-nuclear accumulation energetically favourable. Spelling out the alternative is the cleanest way to see this. \emph{If} the densities had to decay to zero between the nuclei --- as they do at infinity for an isolated atom --- accumulating density in the inter-nuclear region would require steep gradients and a balancing kinetic-energy cost that would offset the potential-energy gain. There would be no bond.

The Bernoulli condition of Chapter~\ref{ch:equation} forbids this decay-to-zero scenario at the inter-electronic interface. The condition states that at the free boundary $\Gamma_{ij}$ between two adjacent electron territories:
\begin{enumerate}
\item $\psi_i$ is continuous from within $\Omega_i$ toward the interface.
\item $\partial \psi_i / \partial n = 0$ on $\Gamma_{ij}$ (homogeneous Neumann).
\item $\psi_i \ne 0$ on the $\Omega_i$ side of the interface.
\end{enumerate}
The first and third together mean that each density is non-zero up to the interface. The second means that the wave function approaches the interface flat --- not steeply --- so that the kinetic-energy cost of the meeting is small. The result is that the inter-nuclear region can accumulate electron density at low kinetic-energy cost while the kernel-attraction potential energy is enhanced by the proximity of that density to both nuclei.

\textbf{Covalent bonding in RealQM is therefore a direct consequence of the free-boundary structure}: meeting non-zero densities allow the kernel potential energy to decrease through inter-nuclear accumulation without a balancing increase in kinetic energy. This is the bonding mechanism, stated in real-space terms.

\paragraph{Connection to R\"udenberg--Kutzelnigg.}
The conclusion is consistent with the R\"udenberg--Kutzelnigg analysis of the chemical bond~\cite{Ruedenberg62, Kutzelnigg90}, in which a careful decomposition of the Slater-determinant wave function shows that covalent bonding is primarily a \emph{kinetic-energy lowering} driven by orbital interference. The unit-density formulation reaches the same qualitative conclusion through a different route: the free-boundary partition rather than orbital overlap. In both pictures, the bond is what allows the kinetic-energy density to redistribute --- in RealQM by flattening at the inter-electronic boundary, in the StdQM picture by interference of overlapping orbitals --- with a net reduction in kinetic energy paid for by the increased kernel attraction.

The two pictures agree on the physics and disagree on the mathematical object. R\"udenberg--Kutzelnigg works inside the StdQM framework and recovers the bond mechanism by analysing antisymmetric one-electron orbitals on $\mathbb{R}^3$ with overlap. RealQM works in a different framework --- non-overlapping orbitals with a free boundary --- and recovers the same mechanism by analysing what happens at the boundary. The convergence of the two routes on the same physical conclusion is a non-trivial cross-check on the multiphase formulation.

\section{The Helium limit: $R \to 0$ with kernels merged}\label{sec:he-limit}

A revealing limit is $R \to 0$. As the two H nuclei approach, their nuclear--nuclear repulsion $1/R$ diverges, so H$_2$ has a finite equilibrium separation. But if the two H$^+$ point kernels are \emph{merged} into a single $+2$ point kernel --- equivalently $R = 0$ with the charges combined --- the system becomes Helium ($Z = 2$, two electrons), and there is no inter-nuclear repulsion to oppose the merger.

The Bernoulli partition is now two unit densities that do not overlap, both bound to the inner shell of the $+2$ kernel, with their interface bisecting the kernel through a plane. This is \textbf{exactly the Level-1 description of the He ground state} from Chapter~\ref{ch:atoms}. The configuration is the He inner shell.

\paragraph{Continuous family from He to two H atoms.}
The H$_2$ bond and the He atom are therefore not separate phenomena. They are the same minimum-energy electron-packing principle applied at different nuclear separations. The continuous family parameterised by $R$ interpolates between:
\begin{itemize}
\item $R = 0$: He atom. Kernels merged into a single $+2$ point kernel. Bernoulli partition bisects the kernel.
\item $R = R_{\rm eq} \approx 1.4$~a.u.: H$_2$ at the variational minimum. Bernoulli partition bisects the inter-nuclear axis. Each electron is bound to both nuclei.
\item $R \to \infty$: two isolated H atoms. Bernoulli partition far from both kernels; each electron is bound to its own H$^+$.
\end{itemize}
Each configuration is a stationary point of the same energy functional. The total energy as a function of $R$ has a minimum at $R_{\rm eq}$ for H$_2$ (the bond energy) and continues, with the nuclear repulsion subtracted, to the He energy at $R = 0$.

The unification clarifies why a single architecture handles both atoms and molecules: \emph{bonding is what the variational principle does when more than one nucleus is present}, with the inter-nuclear repulsion setting the equilibrium separation. There is no separate ``molecule'' machinery; molecules are atoms whose nuclei sit close enough to share electron territories.

\section{What the H$_2$ result establishes}\label{sec:h2-meaning}

The parameter-free Level-1 H$_2$ result fixes several things that the rest of Part~III depends on.

\paragraph{Variational adequacy at Level 1.} The covalent bond is recovered at parameter-free Level 1 with the total energy over-bound by $8.6\%$, the bond length short by $2.3\%$, and the force constant bracketing the exact value --- the same accuracy the architecture delivers for atoms, by the same mechanism and with the same error. This rules out the concern that Level 1 might be adequate for atoms but require extension for bonding. It is not: the architecture handles both, and handles them equally well, which is the relevant point.

\paragraph{The Bernoulli condition does the bonding work.} The mechanism is the non-zero meeting of densities at the free boundary, made energetically favourable by the homogeneous Neumann condition. The bond is not a separately added ingredient; it follows from the free-boundary structure of the multiphase formulation. The same Bernoulli condition that closes the equation in Chapter~\ref{ch:equation} is what makes covalent bonding possible.

\paragraph{Atoms and molecules are the same problem.} The continuous interpolation from He ($R = 0$) to H$_2$ ($R = R_{\rm eq}$) to two H atoms ($R \to \infty$) shows that the multiphase variational problem is one problem parametrised by nuclear positions, not two separate problems (atoms and molecules) handled by different machinery. The hierarchy of reductions (Chapter~\ref{ch:hierarchy}) is therefore an organisational convenience for managing computational cost; it is not a reflection of a physical separation between atomic and molecular regimes.

\paragraph{What is not yet established.} The H$_2$ result is a single bond between identical light atoms with no inner shells. It does not address: bonds between unequal partners (heteronuclear bonds), bonds with inner-shell screening (heavier atoms), multiple bonds (the carbon double bond), aromatic systems, or non-covalent interactions (H-bonds, dispersion). The following chapters take these up in turn, working through Level 3 architectures with kernel parameters $(Z_{\rm kernel}, r_c)$ calibrated against Level 1 for the relevant valence sector.

\section{The H$_2$ result as a foundation}\label{sec:h2-foundation}

A theory of bonding is worth taking seriously to the extent that its simplest case is solid, and to the extent that the case is honestly reported. The H$_2$ result is the simplest case: two atoms, one bond, no inner shells, no fitted parameters, a total energy over-bound by $8.6\%$, a bond length short by $2.3\%$, and an exact energy bracketed by the model's own two constructions at $N=2$. Everything in Chapters~\ref{ch:hydrides} through~\ref{ch:condensed} is built on this foundation. If the H$_2$ result did not survive a parameter-free Level-1 test, the higher-level architectures would have no claim to be calibrated reductions of the same equation. They are calibrated reductions; the H$_2$ result is the evidence that the underlying equation is right.

% --- End of Chapter 7 ---
